\section{Turn each value into a character}
The alphabet has four ordinary ranges and two final symbols. Keep the numerical
sextet separate from the ASCII code stored in output: decimal 19 is an index, while
ASCII \texttt{T} has decimal code 84.
\tagpdfsetup{table/header-rows={1}}
\begin{center}
\begin{tabular}{lll}
Sextet values & Characters & ASCII calculation \\
\hline
0--25 & \texttt{A} through \texttt{Z} & \(q+65\) \\
26--51 & \texttt{a} through \texttt{z} & \(q+71\) \\
52--61 & \texttt{0} through \texttt{9} & \(q-4\) \\
62 & \texttt{+} & 43 \\
63 & \texttt{/} & 47 \\
\end{tabular}
\end{center}

The scalar encoder uses a 64-byte alphabet table. Add the sextet to the table's full
address, load the selected unsigned byte, then store it at the corresponding output
position. A one-byte table entry needs no scale factor. With \texttt{t0}, \texttt{t1}
and \texttt{t2} holding the three input bytes and \texttt{s4} the alphabet pointer,
the first two output positions are:

\Needspace{16\baselineskip}
\MARSCodeLabel{Keep sextet formation separate from the alphabet lookup.}
\begin{minted}[samepage=true]{mars-idea}
    srl $t3, $t0, 2
    daddu $t4, $s4, $t3
    lbu $t4, 0($t4)
    sb $t4, 0($s2)
    andi $t3, $t0, 3
    sll $t3, $t3, 4
    srl $t4, $t1, 4
    or $t3, $t3, $t4
    daddu $t4, $s4, $t3
    lbu $t4, 0($t4)
    sb $t4, 1($s2)
\end{minted}

Use \texttt{lbu} for both input and table bytes. A signed byte load turns a binary
input such as hexadecimal \texttt{ff} into a negative register value, which can change
a shift or produce an out-of-range table address. The table pointer still needs
\texttt{daddu}; the fact that the selected entry occupies one byte does not reduce its
address to a word.
